% GENERATED FILE. Edit canonical lessons, then run npm run generate:books.
% canonical-source-hash: fnv1a64:a5deecbf2a621836
% canonical-lessons: MR-C226-sucavne, MR-C226-parat-dene, MR-C226-nakarne, MR-C226-dhoka-patkarne, MR-C226-sodavne

\chapter{To suggest, To give back, To refuse, To take a risk, To solve}
\label{ch:mr-to-suggest-to-give-back-to-refuse-to-take-a-risk-to-solve}
\hlchaptermodality{226}

\begin{chapteropening}
\textbf{By the end of this chapter:} \emph{I can say to suggest, to give back, to refuse, to take a risk and to solve.}

Complete the last lesson of chapter 226: I can say to suggest, to give back, to refuse, to take a risk and to solve.
\end{chapteropening}

\section[\texorpdfstring{\mr{सुचवणे} (sucavṇe)}{सुचवणे (sucavṇe)}]{\mr{सुचवणे} (sucavṇe) — to suggest}

\label{lesson:MR-C226-sucavne}



\begin{quote}
Before the new one: say the Marathi for to depend, then the Marathi for to avoid.
\end{quote}

\subsection*{You'll want to know: \mr{सुचवणे}}
\textbf{\mr{सुचवणे}} — \emph{sucavṇe} — \textquotedblleft{}to suggest\textquotedblright{}.

\begin{grammarlens}[title={What you've built}]
One more word for this chapter.
\end{grammarlens}

\subsection*{Your turn}
\begin{itemize}
\raggedright
  \item \emph{Say it:} \emph{sucavṇe}
  \item \emph{Say it:} \emph{sucavṇe}, once more
  \item \emph{Say it:} say \emph{sucavṇe}
  \item \emph{Recall:} say the Marathi for to depend, then the Marathi for to avoid, then say \emph{sucavṇe} again
\end{itemize}

\begin{tcolorbox}[breakable,colback=teal!4,colframe=teal!35!black,title={Before you move on}]
What does \mr{सुचवणे} mean? (\textquotedblleft{}To suggest\textquotedblright{}.) Say it once more.
\end{tcolorbox}
\section[\texorpdfstring{\mr{परत} \mr{देणे} (parat deṇe)}{परत देणे (parat deṇe)}]{\mr{परत} \mr{देणे} (parat deṇe) — to give back}

\label{lesson:MR-C226-parat-dene}



\begin{quote}
Before the new one: say the Marathi for to avoid, then the Marathi for to suggest.
\end{quote}

\subsection*{You'll want to know: \mr{परत} \mr{देणे}}
\textbf{\mr{परत} \mr{देणे}} — \emph{parat deṇe} — \textquotedblleft{}to give back\textquotedblright{}.

\begin{grammarlens}[title={What you've built}]
One more word for this chapter.
\end{grammarlens}

\subsection*{Your turn}
\begin{itemize}
\raggedright
  \item \emph{Say it:} \emph{parat deṇe}
  \item \emph{Say it:} \emph{parat deṇe}, once more
  \item \emph{Say it:} say \emph{parat deṇe}
  \item \emph{Recall:} say the Marathi for to avoid, then the Marathi for to suggest, then say \emph{parat deṇe} again
\end{itemize}

\begin{tcolorbox}[breakable,colback=teal!4,colframe=teal!35!black,title={Before you move on}]
What does \mr{परत} \mr{देणे} mean? (\textquotedblleft{}To give back\textquotedblright{}.) Say it once more.
\end{tcolorbox}
\section[\texorpdfstring{\mr{नाकारणे} (nākārṇe)}{नाकारणे (nākārṇe)}]{\mr{नाकारणे} (nākārṇe) — to refuse}

\label{lesson:MR-C226-nakarne}



\begin{quote}
Before the new one: say the Marathi for to suggest, then the Marathi for to give back.
\end{quote}

\subsection*{You'll want to know: \mr{नाकारणे}}
\textbf{\mr{नाकारणे}} — \emph{nākārṇe} — \textquotedblleft{}to refuse\textquotedblright{}.

\begin{grammarlens}[title={What you've built}]
One more word for this chapter.
\end{grammarlens}

\subsection*{Your turn}
\begin{itemize}
\raggedright
  \item \emph{Say it:} \emph{nākārṇe}
  \item \emph{Say it:} \emph{nākārṇe}, once more
  \item \emph{Say it:} say \emph{nākārṇe}
  \item \emph{Recall:} say the Marathi for to suggest, then the Marathi for to give back, then say \emph{nākārṇe} again
\end{itemize}

\begin{tcolorbox}[breakable,colback=teal!4,colframe=teal!35!black,title={Before you move on}]
What does \mr{नाकारणे} mean? (\textquotedblleft{}To refuse\textquotedblright{}.) Say it once more.
\end{tcolorbox}
\section[\texorpdfstring{\mr{धोका} \mr{पत्करणे} (dhokā patkarṇe)}{धोका पत्करणे (dhokā patkarṇe)}]{\mr{धोका} \mr{पत्करणे} (dhokā patkarṇe) — to take a risk}

\label{lesson:MR-C226-dhoka-patkarne}



\begin{quote}
Before the new one: say the Marathi for to give back, then the Marathi for to refuse.
\end{quote}

\subsection*{You'll want to know: \mr{धोका} \mr{पत्करणे}}
\textbf{\mr{धोका} \mr{पत्करणे}} — \emph{dhokā patkarṇe} — \textquotedblleft{}to take a risk\textquotedblright{}.

\begin{grammarlens}[title={What you've built}]
One more word for this chapter.
\end{grammarlens}

\subsection*{Your turn}
\begin{itemize}
\raggedright
  \item \emph{Say it:} \emph{dhokā patkarṇe}
  \item \emph{Say it:} \emph{dhokā patkarṇe}, once more
  \item \emph{Say it:} say \emph{dhokā patkarṇe}
  \item \emph{Recall:} say the Marathi for to give back, then the Marathi for to refuse, then say \emph{dhokā patkarṇe} again
\end{itemize}

\begin{tcolorbox}[breakable,colback=teal!4,colframe=teal!35!black,title={Before you move on}]
What does \mr{धोका} \mr{पत्करणे} mean? (\textquotedblleft{}To take a risk\textquotedblright{}.) Say it once more.
\end{tcolorbox}
\section[\texorpdfstring{\mr{सोडवणे} (soḍavṇe)}{सोडवणे (soḍavṇe)}]{\mr{सोडवणे} (soḍavṇe) — to solve}

\label{lesson:MR-C226-sodavne}



\begin{quote}
Before the new one: say the Marathi for to refuse, then the Marathi for to take a risk.
\end{quote}

\subsection*{You'll want to know: \mr{सोडवणे}}
\textbf{\mr{सोडवणे}} — \emph{soḍavṇe} — \textquotedblleft{}to solve\textquotedblright{}.

\begin{grammarlens}[title={What you've built}]
One more word for this chapter.
\end{grammarlens}

\subsection*{Your turn}
\begin{itemize}
\raggedright
  \item \emph{Say it:} \emph{soḍavṇe}
  \item \emph{Say it:} \emph{soḍavṇe}, once more
  \item \emph{Say it:} say \emph{soḍavṇe}
  \item \emph{Recall:} say the Marathi for to refuse, then the Marathi for to take a risk, then say \emph{soḍavṇe} again
\end{itemize}

\begin{tcolorbox}[breakable,colback=teal!4,colframe=teal!35!black,title={Before you move on}]
What does \mr{सोडवणे} mean? (\textquotedblleft{}To solve\textquotedblright{}.) Say it once more.
\end{tcolorbox}
